% !TeX program = xelatex
% Сборка: два запуска xelatex TiAl_solvus_1000K.tex в одной папке.
% Оба векторных рисунка и библиография включены непосредственно в исходник.
\documentclass[11pt,a4paper]{article}
\usepackage{fontspec}
\usepackage[russian]{babel}
\IfFontExistsTF{Cambria}{\setmainfont{Cambria}}{\setmainfont{DejaVu Serif}}
\IfFontExistsTF{Calibri}{\setsansfont{Calibri}}{\setsansfont{DejaVu Sans}}
\IfFontExistsTF{Consolas}{\setmonofont[Scale=MatchLowercase]{Consolas}}{\setmonofont{DejaVu Sans Mono}}
\usepackage[a4paper,left=23mm,right=21mm,top=21mm,bottom=20mm,headheight=15pt,headsep=7mm,footskip=10mm]{geometry}
\usepackage{amsmath,amssymb,booktabs,array,xcolor,titlesec,fancyhdr,caption,pgfplots,hyperref}
\pgfplotsset{compat=1.16}
\definecolor{ink}{HTML}{213542}
\definecolor{teal}{HTML}{006D70}
\definecolor{rust}{HTML}{A8492B}
\definecolor{muted}{HTML}{526673}
\definecolor{linegray}{HTML}{C5D0D3}
\hypersetup{unicode=true,colorlinks=true,linkcolor=teal,citecolor=teal,urlcolor=teal,
 pdftitle={Равновесные составы фаз α и α₂ в Ti–Al при 1000 К},pdfsubject={Кластерное разложение, свободная энергия и прямой расчёт на MTP}}
\urlstyle{same}
\captionsetup{font=small,labelfont={bf,color=teal},labelsep=period,singlelinecheck=false}
\titleformat{\section}{\sffamily\bfseries\Large\color{ink}}{\color{teal}\thesection}{.55em}{}
\titleformat{\subsection}{\sffamily\bfseries\large\color{ink}}{\thesubsection}{.55em}{}
\titlespacing*{\section}{0pt}{16pt plus 3pt}{6pt}
\titlespacing*{\subsection}{0pt}{10pt plus 2pt}{4pt}
\setlength{\parindent}{1.1em}
\setlength{\parskip}{2.5pt plus .7pt}
\setlength{\textfloatsep}{11pt plus 2pt minus 2pt}
\setlength{\intextsep}{10pt plus 2pt minus 2pt}
\setlength{\floatsep}{8pt plus 2pt}
\setlength{\emergencystretch}{2em}
\brokenpenalty=10000 \widowpenalty=8000 \clubpenalty=8000
\renewcommand{\arraystretch}{1.14}
\raggedbottom
\pagestyle{fancy}\fancyhf{}
\fancyhead[L]{\small\sffamily\color{muted}Ti--Al\enspace/\enspace Свободная энергия и фазовое равновесие}
\fancyhead[R]{\small\sffamily\color{muted}$1000$ К}
\fancyfoot[R]{\small\sffamily\color{muted}\thepage}
\renewcommand{\headrulewidth}{0.25pt}
\begin{document}
\thispagestyle{plain}
{\sffamily\small\color{teal}\bfseries НАУЧНЫЙ ОТЧЁТ\hfill 25 ноября 2026 г.\par}
\vspace{4mm}
{\sffamily\bfseries\fontsize{24}{28}\selectfont\color{ink}
Равновесные составы фаз \(\boldsymbol{\alpha}\) и \(\boldsymbol{\alpha_2}\)\\
в Ti--Al при 1000 К\par}
\vspace{2mm}
{\sffamily\large\color{muted}Прямая атомистическая оценка при нулевом давлении\\
и сравнение с кластерным разложением C46\par}
\vspace{4mm}\noindent{\color{teal}\rule{\linewidth}{1.1pt}}\par\vspace{1mm}

\noindent При свободном объёме и отношении осей \(c/a\) расчёт на потенциале
Qi и соавторов дал \textbf{14,98 ат.\% Al в \(\boldsymbol{\alpha}\)} и
\textbf{22,44 ат.\% Al в \(\boldsymbol{\alpha_2}\)}. Условные 95\%-е интервалы
составили \(14{,}92\)--\(15{,}06\) и \(22{,}33\)--\(22{,}55\) ат.\%.
Это равновесие двух кристаллических фаз в явно заданной области тепловых
смещений. Границы проверены на четырёх размерах, вплоть до 1400 атомов.
При изменении допустимых смещений и независимом повторении расчёта
измеренные составы менялись на величины порядка \(0{,}1\)--\(0{,}2\) процентного пункта.

\begin{center}
\begin{tikzpicture}
\begin{axis}[width=.91\linewidth,height=47mm,axis x line*=bottom,axis y line=none,
 xmin=14.0,xmax=23.15,ymin=.45,ymax=3.53,
 xtick={14,15,16,17,18,19,20,21,22,23},enlarge x limits=false,
 tick label style={font=\sffamily\footnotesize,text=muted},
 tick style={draw=muted},xlabel={Al, ат.\%},xlabel style={font=\sffamily\small},
 xmajorgrids,grid style={linegray!65},clip=false]
\addplot[teal,line width=1.5pt,mark=*,mark size=2.3pt] coordinates {(14.982,3) (22.445,3)};
\addplot[rust,line width=1.5pt,mark=square*,mark size=2.2pt] coordinates {(16.083,2) (22.159,2)};
\addplot[muted,line width=1.25pt,dashed,mark=triangle*,mark size=2.5pt] coordinates {(14.958,1) (21.958,1)};
\node[fill=white,text=teal,font=\sffamily\footnotesize,inner xsep=5pt] at (axis cs:18.60,3) {MTP, свободное \(c/a\)};
\node[fill=white,text=rust,font=\sffamily\footnotesize,inner xsep=5pt] at (axis cs:18.60,2) {MTP, \(c/a\) как у C46};
\node[fill=white,text=muted,font=\sffamily\footnotesize,inner xsep=5pt] at (axis cs:18.60,1) {C46, фиксированное \(c/a\)};
\node[above,text=ink,font=\small] at (axis cs:14.982,3.09) {\(\alpha\)};
\node[above,text=ink,font=\small] at (axis cs:22.445,3.09) {\(\alpha_2\)};
\end{axis}
\end{tikzpicture}
\captionof{figure}{Интервалы сосуществования при \(p=0\). Концы каждого
отрезка соответствуют равновесным составам двух фаз одной модели.
У прямого MTP \(N=1400\) и \(W<0{,}40\), у C46 \(N=5120\).}
\label{fig:tielines}
\end{center}

\noindent\textbf{Влияние формы ячейки.}
При одинаковом \(c/a\) прямой MTP даёт на \(1{,}12\) п.п. больше Al
в \(\alpha\) и на \(0{,}20\) п.п. больше в \(\alpha_2\), чем C46.
Освобождение отношения осей в атомистическом расчёте сдвигает границы
на \(-1{,}10\) и \(+0{,}29\) п.п. Таким образом, совпадение результатов
для \(\alpha\) при сравнении \(14{,}98\) с \(14{,}96\) возникает из-за
компенсации двух эффектов.

\noindent\textbf{Следующий расчёт.}
При тех же \(T,p,c/a\) целесообразно построить классический гармонический
вариант C46 и повторить определение обеих границ. Это даст явную оценку
влияния статистики ядер. Для измерения ошибки самого потенциала
нужна независимая разность свободных энергий MTP--DFT для обеих фаз.

\vfill
\noindent{\small\sffamily\color{muted}
Малые ячейки и ангармонизм — с.~\pageref{sec:smallcells}.
Атомистический ансамбль — с.~\pageref{sec:ensemble}.
Составы и сходимость — с.~\pageref{sec:bulk-results}.
Дальнейшие расчёты — с.~\pageref{sec:outlook}.}
\clearpage
\section{Постановка задачи}
На богатой титаном стороне диаграммы Ti--Al фаза \(\alpha\) представляет
собой замещающий раствор Al в гексагональной плотноупакованной (ГПУ)
решётке. В фазе \(\alpha_2\) структуры D0$_{19}$ алюминий преимущественно
занимает одну из четырёх подрешёток. В идеальном \(\mathrm{Ti_3Al}\)
её заполнение полное, тогда как равновесная при конечной температуре
фаза допускает отклонение от стехиометрии.

Обозначим атомную долю алюминия через \(x=N_{\mathrm{Al}}/N\), а
свободную энергию Гиббса на атом фазы \(i\) через \(g_i(x,T,p)\).
Концы двухфазного интервала удовлетворяют условию общей касательной,
\begin{align}
  g'_\alpha(x_\alpha)&=g'_{\alpha_2}(x_{\alpha_2})=\Delta\mu, \\
  g_\alpha(x_\alpha)-x_\alpha\Delta\mu
    &=g_{\alpha_2}(x_{\alpha_2})-x_{\alpha_2}\Delta\mu ,
  \qquad \Delta\mu=\mu_{\rm Al}-\mu_{\rm Ti}.
\end{align}
Малый систематический сдвиг разности свободных энергий способен заметно
сместить касательную. Изменение свободной энергии одной декорации
нужно отличать от сдвига этой касательной после усреднения по всем
химическим заполнениям и атомным смещениям.

В проектной модели C46 энергии и колебательные вклады представлены
кластерным разложением. Исходные атомистические данные получены на
потенциале моментных тензоров (MTP) Qi и соавт.~\cite{qi2023}.
Воспроизведённая версия C46 использует авторские коэффициенты из комплекта
рукописи v16. Независимая подгонка на собственной выборке структур давала
при 1000 К \(14{,}764/22{,}707\) ат.\% Al. Сравнения ниже выполнены с
авторским набором \(14{,}958/21{,}958\) ат.\% Al.
В расчёте C46 химическое разупорядочение учитывалось методом
Монте-Карло, а колебания — в квантовом квазигармоническом приближении.
В настоящей работе использован тот же файл потенциала. Сначала фиксированные
декорации и полный перебор малых ячеек позволили проверить отдельные
вклады в свободную энергию. Затем для макроскопических составов мы
восстановили равновесные веса фаз из атомистических ансамблей без
кластерной аппроксимации энергии. С экспериментом прямые свободные
энергии Гиббса в этом расчёте не калибровались.

\section{Проверка свободных энергий в малых ячейках}
\label{sec:smallcells}

Малый размер ячейки позволил выполнить полный перебор химических заполнений и
на одних и тех же конфигурациях измерить два эффекта. Прямые фононы показывают,
насколько кластерное разложение C46 воспроизводит гармонические свободные
энергии MTP. Переход от гармонического гамильтониана к полному MTP при
классической статистике ядер измеряет вклад ангармонизма. Все расчёты этого
раздела относятся к $T=1000$~К, фиксированным $V/N=17{,}28833$~\AA$^3$/атом
и $c/a=4{,}6046/2{,}9446$. Здесь вычислена свободная энергия Гельмгольца.

\subsection{Интегрирование от гармонической опоры}

Для химического заполнения $\sigma$ сначала релаксировали внутренние
координаты и рассчитывали гессиан $H_\sigma$ того же MTP. После исключения трёх
общих переносов гармоническая энергия равна
$U_{\mathrm h}=U_{\min}+\tfrac12\mathbf z^{\mathsf T}H_\sigma\mathbf z$,
где $\mathbf z$ отсчитан от релаксированного минимума. При одной и той же
области допустимых смещений $D$ во всех промежуточных состояниях используем
\begin{align}
 U_\lambda &= U_{\mathrm h}+\lambda(U_{\mathrm{MTP}}-U_{\mathrm h}),\\
 F_{\sigma,\mathrm{MTP}}^{D}-F_{\sigma,\mathrm h}^{D}
 &=\int_0^1
     \langle U_{\mathrm{MTP}}-U_{\mathrm h}\rangle_{\lambda,\sigma,D}
     \,\mathrm d\lambda .
\end{align}
В восьмиатомном пилоте средние вычисляли перевзвешиванием независимых пробных
смещений. Интеграл на сетке из 65 значений $\lambda$ проверяли по более редкой
сетке и по отношению статистических сумм в конечных точках. Для количественной
проверки ограниченных свободных энергий отдельно оценивали долю гармонического
распределения, которая попала в $D$. При $k$ атомах Al итоговая каноническая
энергия ячейки учитывает кратность $g_\sigma$ каждого класса эквивалентных
заполнений,
\[
 F_k^D=-k_{\mathrm B}T\ln
   \sum_{\sigma\,:\,N_{\mathrm{Al}}=k}g_\sigma
     \exp[-F_\sigma^D/(k_{\mathrm B}T)] .
\]

\subsection{Короткие расстояния и выбор кристаллического бассейна}

В независимой проверке межслойного заполнения $\mathrm{Ti}_6\mathrm{Al}_2$
сгенерировали 65\,536 пробных смещений из смеси гауссовского распределения
и многомерного распределения Стьюдента. В одном состоянии два атома Ti
сблизились до $0{,}636816$~\AA\ при сохранении каждого атома в своей исходной
ячейке Вороного. Полная энергия восьми атомов упала с $-56{,}069$ до
$-128\,392{,}623$~эВ. Почти весь статистический вес пришёлся на один снимок.
Повторный расчёт MLIP и периодическое размножение до 64 атомов воспроизвели
аномалию. Расхождение энергий на атом при размножении составило
$2{,}85\cdot10^{-9}$~эВ. На прямолинейном пути из 101 точки между релаксированной
структурой и этим снимком энергия сначала поднималась на $42{,}20$~эВ/ячейку
над энергией минимума. Высота относится к конкретному пути и объясняет, почему
локальные траектории способны долго удерживаться вблизи кристалла.

Кристаллический бассейн задали через смещение каждого атома от идеального
ГПУ-узла после удаления общего переноса. Для каждой грани его ячейки
Вигнера--Зейтца вычисляли нормальную проекцию смещения в долях расстояния
от узла до грани. Максимум этих безразмерных величин по всем атомам и
граням обозначили $W$. При $W<0{,}40$ минимально возможное расстояние в
этой ячейке ограничено снизу величиной
$(1-0{,}40)\cdot 2{,}861976 = 1{,}71719$~\AA. Параметр радиального базиса
MTP равен $1{,}68052$~\AA. Опубликованный файл потенциала допускает
вычисление энергии и на меньших расстояниях. Геометрическое ограничение
поэтому входит в определение измеренного кристаллического равновесия.

\subsection{Предпочтение порядка \texorpdfstring{D0$_{19}$}{D0₁₉} в ячейке из восьми атомов}

При $N=8$ и $k=2$ имеется $\binom82=28$ заполнений. Группировка по симметрии
даёт 4 варианта D0$_{19}$, 12 вариантов с базальной парой Al и 12 вариантов
с межслойной парой другого типа. В табл.~\ref{tab:small8} приведена энергия
одного представителя каждого класса относительно D0$_{19}$. Кратности входят
в статистическую сумму только при последующем объединении заполнений.

\begin{table}[htbp]
\centering\small
\caption{Энергетическая цена двух заполнений $\mathrm{Ti}_6\mathrm{Al}_2$
относительно D0$_{19}$, мэВ/атом. Знак $\pm$ здесь и в
табл.~\ref{tab:small16} обозначает одну стандартную ошибку.}
\label{tab:small8}
\begin{tabular}{@{}lrr@{}}
\toprule
Модель & Базальная пара & Межслойная пара\\
\midrule
C46, статическая энергия & $78{,}60$ & $110{,}61$\\
Прямой MTP, статическая энергия & $84{,}10$ & $112{,}92$\\
C46, квантовая гармоника & $61{,}01$ & $85{,}66$\\
Прямой MTP, квантовая гармоника & $65{,}43$ & $86{,}44$\\
Классическая гармоника, $W<0{,}40$ & $65{,}55\pm0{,}05$ & $87{,}78\pm0{,}05$\\
Классический MTP, $W<0{,}40$ & $68{,}07\pm0{,}16$ & $89{,}98\pm0{,}10$\\
\bottomrule
\end{tabular}
\end{table}

Для прямой квантовой гармоники силовые постоянные строили в суперячейках
$5\times5\times6$ исходных восьмиатомных ячеек (1200 атомов), фононную
сумму вычисляли на сетке $16^3$ в зоне Бриллюэна. При общей гармонической
модели C46 занижает цену базальной пары на $4{,}42$~мэВ/атом, межслойной
на $0{,}78$~мэВ/атом. В восьмиатомной классической ячейке ангармонизм
повышает эти цены соответственно на $2{,}52\pm0{,}13$ и
$2{,}20\pm0{,}08$~мэВ/атом по согласованным разностям. Независимая широкая
выборка для всех трёх заполнений дала $2{,}08\pm0{,}17$ и
$1{,}60\pm0{,}65$~мэВ/атом. Обе серии дают дополнительную стабилизацию порядка.
Для базальной пары расхождение оценок составляет около двух объединённых
стандартных ошибок. Квантовые результаты включают колебания во всей зоне
Бриллюэна, классические относятся к модам малой периодической ячейки.
Суммирование двух измеренных сдвигов потребовало бы согласования этих условий.

\subsection{Нелинейность свободной энергии по составу в ячейке на 16 атомов}

Для $N=16$ и $k=2,3,4$ атомные доли Al равны $12{,}5$, $18{,}75$ и
$25$~ат.\,\%. Перечислены соответственно 120, 560 и 1820 заполнений. После
учёта симметрии остались 43 различных конфигурации с устойчивыми гессианами
по 45 внутренним степеням свободы. Меру нелинейности определили как
\[
 \delta f_{16}=\frac{F_3-(F_2+F_4)/2}{16}.
\]
Положительная $\delta f_{16}$ означает, что при среднем составе свободная
энергия лежит выше прямой между крайними составами. Комбинация сокращает
общие постоянные и линейные по $k$ члены, включая классический вклад масс
от интегрирования по импульсам. У статических гамильтонианов после
суммирования химических заполнений $\delta f_{16}$ составила
$3{,}015$~мэВ/атом для C46 и $3{,}706$~мэВ/атом для прямого MTP.

\begin{table}[htbp]
\centering\small
\caption{Значение $\delta f_{16}$ и его ангармоническая поправка, мэВ/атом.
Столбцы представляют независимые серии. Во второй использована смесь
гауссовского распределения и распределения Стьюдента с пятью степенями свободы.}
\label{tab:small16}
\begin{tabular}{@{}lrr@{}}
\toprule
Величина & Смесь гауссиан & Смесь со Стьюдентом\\
\midrule
Гармоника, $W<0{,}40$ & $1{,}080\pm0{,}019$ & $1{,}077\pm0{,}009$\\
Полный MTP, $W<0{,}40$ & $1{,}374\pm0{,}085$ & $1{,}383\pm0{,}036$\\
Ангармоническая поправка, $W<0{,}40$ & $+0{,}294\pm0{,}071$ & $+0{,}305\pm0{,}032$\\
Ангармоническая поправка, $W<0{,}35$ & $+0{,}291\pm0{,}069$ & $+0{,}307\pm0{,}028$\\
\bottomrule
\end{tabular}
\end{table}

Первые оценки ангармонической поправки, $+0{,}640\pm0{,}090$ и
$-0{,}233\pm0{,}263$~мэВ/атом, расходились даже по знаку. Разбор вкладов
выявил четыре класса заполнений с большим статистическим весом. Для каждого
добавили по 32\,768 новых проб в каждой из двух серий. Для оставшихся
39 классов сохранили соответствующие независимые исходные выборки. После
уточнения обе оценки дают около $+0{,}30$~мэВ/атом, и этот сдвиг сохраняется
при сужении допустимой области (табл.~\ref{tab:small16}). Он характеризует
форму $F(x)$ для трёх заданных составов в ячейке из 16 атомов. Непрерывное
положение двух контактов при $p=0$ определяет отдельный расчёт с флуктуациями
состава, тепловых смещений и объёма.

\section{Прямой атомистический расчёт сольвуса}
\label{sec:ensemble}

\subsection{Степени свободы и механические условия}
В больших ячейках одновременно менялись химические типы узлов, координаты
атомов и объём. В основных сериях независимо менялось отношение осей
\(\eta=c/a\) при гексагональной симметрии. Для контроля геометрии C46
объём также релаксировал, а \(\eta\) фиксировалось. В обоих случаях
во всех вероятностях принятия использовано одно внешнее давление \(p=0\).
Две фазы могут иметь разные равновесные объёмы даже в отсутствие давления.

Чтобы область допустимых смещений двигалась вместе с кристаллом,
для каждого узла взяли вектор смещения \(\mathbf u_i\) относительно
идеального положения в текущей ячейке и вычли среднее
\(\overline{\mathbf u}\). Для двенадцати векторов \(\mathbf d_{ij}\)
от узла к его ближайшим соседям вычисляли
\[
 W=\max_{i,j}\frac{2(\mathbf u_i-\overline{\mathbf u})\cdot\mathbf d_{ij}}
                              {|\mathbf d_{ij}|^2}.
\]
Ход допустим при \(W<W_{\max}\). Условие сохраняли для обоих химических
типов и всех окон. Предложения объёма и формы дополнительно удовлетворяли
\(15<V/N<21\)~\AA$^3$/атом, \(1{,}45<\eta<1{,}72\) при свободных осях
и \(d_{\rm nn}(1-W_{\max})\ge 1{,}68052\)~\AA, где \(d_{\rm nn}\) —
кратчайшее расстояние между идеальными узлами текущей ячейки.
Так сохранялась область, для которой проведена проверка коротких расстояний.
В основной ячейке с \(W_{\max}=0{,}40\) у обеих фаз доля кадров в
пределах \(0{,}25\)~\AA$^3$/атом от границ объёма и \(0{,}01\) от
границ \(\eta\) была нулевой.

\subsection{Химический порядок и зонтичные окна}
Пусть \(n_j\) — число атомов Al на подрешётке \(j\), а \(k=\sum_{j=1}^4 n_j\).
Параметр порядка и энергия зонтика в окне \(l\) равны
\begin{equation}
 Q=\frac{4(4\max_j n_j-k)}{3N},\qquad
 B_l(Q)=\tfrac12 K_l(Q-q_l)^2 .
\end{equation}
У идеального D0$_{19}$ получаем \(Q=1\). В случайном растворе максимум
из четырёх заполнений даёт положительный \(Q\) при конечном \(N\).
Определение объединяет четыре трансляционно эквивалентных варианта
химического порядка. Упорядоченные домены разных
вариантов могут снижать общий \(Q\), поэтому проверяли также статистику
состава и распределение заполнений подрешёток.

Для химической замены в случайном узле вероятность принятия определяется
разностью \(U_{\rm MTP}+B_l-\Delta\mu_l k\), где \(\Delta\mu_l\) —
конфигурационное поле, положительное в пользу Al. Каждый проход включает
по \(N\) предложений замены и декартова смещения. Объём предлагается раз
за проход симметрично по \(V\). Для изотропного масштабирования координат
логарифм отношения весов содержит
\[
 -\beta[\Delta U_{\rm MTP}+p\,\Delta V]+N\ln(V'/V),
 \qquad \beta=(k_{\rm B}T)^{-1}.
\]
Относительные координаты после удаления переноса дают степень \(N-1\),
интегрирование общего переноса возвращает ещё множитель \(V\).
Для свободных осей каждые три прохода предлагали симметричное изменение
\(\ln\eta\) при постоянном объёме. В расчёте химических замен использованы
одинаковые фиктивные массы. Для физических классических масс
\(\Delta\mu_{\rm phys}=\Delta\mu_{\rm cfg}
 -\tfrac32 k_{\rm B}T\ln(m_{\rm Al}/m_{\rm Ti})\).
При 1000 К поправка составляет \(+0{,}07410\) эВ на замену.
Линейная по составу поправка к свободной энергии смещает наклон общей
касательной и сохраняет составы в точках касания.

\subsection{Как из окон получены две фазовые границы}
Центры зонтиков охватывали \(Q=0{,}035\)--\(0{,}910\).
Соседние окна перекрываются по заполнениям. При одинаковых \(N\),
\(W_{\max}\) и механических условиях общая атомистическая энергия и
мера координат сокращаются в отношении весов двух окон. Для их объединения
методом MBAR достаточно точно знать \(k,Q\) и добавленный член
\(\beta[B_l(Q)-\Delta\mu_l k]\) в каждом окне~\cite{mbar2008}.
На общей шкале поля вероятность можно перевзвешивать по формуле
\[
 P_{\Delta\mu'}(k,Q)\ \propto\ P_{\Delta\mu}(k,Q)
                       \exp[\beta(\Delta\mu'-\Delta\mu)k] .
\]
Разделитель \(Q^\ddagger\) выбирали в минимуме между модами
при поле, где их высоты равны. С этим разделителем находили \(\Delta\mu^*\)
из условия
\[
 \frac{P_{\Delta\mu^*}(Q>Q^\ddagger)}
      {P_{\Delta\mu^*}(Q\le Q^\ddagger)}=4 .
\]
Четвёрка задаёт одинаковый термодинамический вес для неупорядоченной фазы
и каждого из четырёх вариантов D0$_{19}$~\cite{janke1993}.
При найденном поле левый и правый условные средние \(\langle k/N\rangle\)
дают составы \(\alpha\) и \(\alpha_2\). Пересечение изотерм
\(\langle x\rangle(\Delta\mu)\) разных размеров служит независимой по
критерию проверкой положения перехода~\cite{borgs1992}.

Статистические интервалы рассчитаны ресэмплированием последовательных
неперекрывающихся блоков внутри каждой траектории. Длину блока задавала
наиболее медленная из величин \(k,Q,V,c/a\), а в каждой бутстрэп-выборке
заново оценивали веса, разделитель и поле. Для \(N=96,288,768,1400\)
использованы соответственно 96, 128, 192 и 192 бутстрэп-реализации.
Эти интервалы характеризуют шум данной выборки. Зависимость от размера,
чувствительность к \(W_{\max}\) и разброс независимых стартов
приведены отдельно.

\section{Равновесные составы при 1000 К}
\label{sec:bulk-results}

При нулевом внешнем давлении, свободных объёме и отношении осей
$c/a$ прямой расчёт на 1400 атомах дал
\[
 x_{\alpha}=14{,}982\ \text{ат.\% Al}, \qquad
 x_{\alpha_2}=22{,}445\ \text{ат.\% Al}.
\]
Для этой оценки выбран более широкий из двух проверенных пределов
смещений, $W<0{,}40$. Условные 95\%-е интервалы по блокам траекторий
равны $14{,}917$--$15{,}063$ и $22{,}327$--$22{,}553$ ат.\% Al.
Ширина интервала сосуществования составляет $7{,}463$ процентного
пункта (п.п.), с 95\%-м интервалом $7{,}323$--$7{,}580$ п.п.
Здесь и далее составы даны в атомных процентах Al, а изменения
составов --- в процентных пунктах.

\subsection{Размер ячейки и статистическое разрешение}

Для каждого из четырёх размеров изучены две области допустимых
смещений. Вся серия охватывает десять физических ансамблей при
различных $N$, пределе $W$ и условиях на $c/a$. В них накоплено
$255\,200$ производственных проходов в 268 зонтичных траекториях.
Ещё 16 независимых траекторий дали $11\,200$ проходов без зонтичного
ограничения. Один проход включает $N$ попыток химической замены
и $N$ попыток локального смещения; изменение объёма предлагается
каждый проход. Первоначальная релаксация исключена из всех средних.

\begin{table}[htbp]
\centering\small
\caption{Составы фаз (ат.\% Al) при свободном $c/a$. Для более
широкой области приведён также барьер $\Delta F_Q$ между максимумами
распределения $Q$, измеренный в поле равных высот.}
\label{tab:bulk-size}
\begin{tabular}{@{}r c rr rr r@{}}
\toprule
 & & \multicolumn{2}{c}{$W<0{,}35$} & \multicolumn{2}{c}{$W<0{,}40$}
 & $W<0{,}40$\\
\cmidrule(lr){3-4}\cmidrule(lr){5-6}
$N$ & Ячейка & $\alpha$ & $\alpha_2$ & $\alpha$ & $\alpha_2$
 & $\Delta F_Q/(k_{\mathrm B}T)$\\
\midrule
  96 & $4\times4\times3$   & $14{,}93$ & $22{,}83$ & $15{,}48$ & $22{,}86$ & $1{,}18$\\
 288 & $6\times6\times4$   & $15{,}04$ & $22{,}47$ & $15{,}19$ & $22{,}42$ & $2{,}17$\\
 768 & $8\times8\times6$   & $15{,}02$ & $22{,}52$ & $14{,}93$ & $22{,}40$ & $4{,}23$\\
1400 & $10\times10\times7$ & $14{,}86$ & $22{,}36$ & $14{,}98$ & $22{,}44$ & $5{,}66$\\
\bottomrule
\end{tabular}
\end{table}

При $W<0{,}40$ увеличение $N$ с 768 до 1400 сдвинуло обе границы
на $+0{,}05$ п.п. Условные 95\%-е интервалы этих разностей равны
$[-0{,}05;\,+0{,}16]$ для $\alpha$ и $[-0{,}11;\,+0{,}21]$ п.п.
для $\alpha_2$. При $W<0{,}35$ на том же шаге получены
$-0{,}166$ и $-0{,}159$ п.п. Для $\alpha$ при более тесном
ограничении интервал $[-0{,}298;\,-0{,}034]$ п.п. указывает на
ещё различимый размерный сдвиг. Поэтому в качестве результата
используется непосредственно измеренная точка при $N=1400$.
По этим размерам и независимым стартам практически достижимое
разрешение сравнения составляет около $0{,}2$ п.п. Бутстрэп
выполнялся по блокам внутри каждой траектории с повторной оценкой
поля и составов; при $N=768$ и 1400 использовано по 192 реализации.
Такой интервал отражает неопределённость выборки внутри заданного
ансамбля. Влияние размера, формы и ограничения смещений оценивается
отдельными расчётами.

При $N=1400$ и $W<0{,}40$ поле сосуществования равно
$\Delta\mu_{\mathrm{cfg}}=3{,}29368$ эВ на замену Ti$\to$Al
в конфигурационной нормировке. Максимумы распределения $Q$ при
этом поле находятся около $0{,}063$ и $0{,}878$. В фазе $\alpha_2$
на наиболее богатой алюминием подрешётке Al занимает в среднем
$87{,}5\%$ узлов; на остальных трёх, упорядоченных по убыванию
заполнения, --- $1{,}19\%$, $0{,}72\%$ и $0{,}37\%$.
Это характерный признак порядка D0$_{19}$ при отклонении от
стехиометрии Ti$_3$Al. Объём на атом и $c/a$ составили
$17{,}457$~\AA$^3$ и $1{,}5915$ для $\alpha$, $17{,}241$~\AA$^3$
и $1{,}6064$ для $\alpha_2$. В пределах выбранных диагностических
полос у границ допустимого объёма и $c/a$ статистический вес обеих
фаз оказался нулевым.

\subsection{Чувствительность к ограничению смещений}

При $N=1400$ расширение области с $W<0{,}35$ до $W<0{,}40$
повысило поле сосуществования на $4{,}59$ мэВ на замену
Ti$\to$Al. Состав $\alpha$ изменился на $+0{,}123$ п.п.
с 95\%-м интервалом от $+0{,}003$ до $+0{,}251$ п.п.
Для $\alpha_2$ сдвиг составил $+0{,}083$ п.п. с интервалом
от $-0{,}079$ до $+0{,}251$ п.п.
Более широкая область принимает дополнительные состояния с
большими тепловыми смещениями. При $N=288$ в ансамбле $W<0{,}40$
вероятность выполнить более строгое условие $W<0{,}35$ составила
$0{,}158$ в $\alpha$ и $0{,}468$ в $\alpha_2$. Поэтому ужесточение
ограничения термодинамически благоприятствует упорядоченной фазе.
Повторное взвешивание отобранных конфигураций дало
$15{,}05$ и $22{,}34$ ат.\% Al. В самостоятельном расчёте
при $W<0{,}35$ получены $15{,}04$ и $22{,}47$ ат.\%.

Условие на $W$ должно выполняться одновременно на всех атомах.
На 1400 атомах вероятность пройти эту же фильтрацию в фазе $\alpha$
упала до $1{,}23\cdot10^{-5}$. Эффективное число сохранившихся
кадров по весам Киша составило всего $3{,}6$ ещё до поправки на
автокорреляцию. Для наибольшей ячейки вывод о чувствительности
основан на самостоятельных сериях при каждой величине $W$.
Обе серии характеризуют конденсированный ГПУ-бассейн данного MTP.


\begin{figure}[htbp]
\centering
\begin{tikzpicture}
\begin{axis}[width=.91\linewidth,height=58mm,xmin=0,xmax=1,ymin=0,
 xlabel={Параметр порядка \(Q\)},ylabel={Плотность \(P(Q)\)},
 tick label style={font=\sffamily\footnotesize}, label style={font=\sffamily\small},
 legend style={font=\sffamily\footnotesize,draw=none,fill=none,
 at={(.05,.97)},anchor=north west,legend columns=2,column sep=5pt},
 axis x line*=bottom,axis y line*=left,axis line style={muted},
 ymajorgrids,y grid style={linegray!65},
 xtick={0,0.2,0.4,0.6,0.8,1.0},xticklabels={0,{0,2},{0,4},{0,6},{0,8},1}]

\addplot[muted, densely dotted,line width=1.1pt,no marks] coordinates {
(0.00250,0.1069668) (0.00750,0.1090155) (0.01250,0.1130348) (0.01750,0.1190158) (0.02250,0.1268636)
(0.02750,0.1364220) (0.03250,0.1476232) (0.03750,0.1602448) (0.04250,0.1741685) (0.04750,0.1892138)
(0.05250,0.2051600) (0.05750,0.2218640) (0.06250,0.2391425) (0.06750,0.2567796) (0.07250,0.2746371)
(0.07750,0.2925419) (0.08250,0.3102976) (0.08750,0.3277742) (0.09250,0.3448139) (0.09750,0.3612269)
(0.10250,0.3769041) (0.10750,0.3916573) (0.11250,0.4053686) (0.11750,0.4179166) (0.12250,0.4291547)
(0.12750,0.4390046) (0.13250,0.4473990) (0.13750,0.4542624) (0.14250,0.4595842) (0.14750,0.4633811)
(0.15250,0.4656604) (0.15750,0.4665045) (0.16250,0.4659781) (0.16750,0.4641936) (0.17250,0.4612873)
(0.17750,0.4573758) (0.18250,0.4526172) (0.18750,0.4471791) (0.19250,0.4411939) (0.19750,0.4348303)
(0.20250,0.4282451) (0.20750,0.4215551) (0.21250,0.4148993) (0.21750,0.4083901) (0.22250,0.4020847)
(0.22750,0.3960773) (0.23250,0.3903675) (0.23750,0.3849923) (0.24250,0.3799433) (0.24750,0.3751598)
(0.25250,0.3706309) (0.25750,0.3663022) (0.26250,0.3620864) (0.26750,0.3579637) (0.27250,0.3538843)
(0.27750,0.3497796) (0.28250,0.3456792) (0.28750,0.3415365) (0.29250,0.3373552) (0.29750,0.3332100)
(0.30250,0.3290780) (0.30750,0.3250506) (0.31250,0.3211735) (0.31750,0.3174700) (0.32250,0.3140446)
(0.32750,0.3109406) (0.33250,0.3081716) (0.33750,0.3058161) (0.34250,0.3038984) (0.34750,0.3023938)
(0.35250,0.3013577) (0.35750,0.3007269) (0.36250,0.3005203) (0.36750,0.3006979) (0.37250,0.3011882)
(0.37750,0.3019816) (0.38250,0.3030239) (0.38750,0.3042319) (0.39250,0.3055899) (0.39750,0.3070456)
(0.40250,0.3085231) (0.40750,0.3100238) (0.41250,0.3114805) (0.41750,0.3128612) (0.42250,0.3141863)
(0.42750,0.3153910) (0.43250,0.3165006) (0.43750,0.3175174) (0.44250,0.3184083) (0.44750,0.3192200)
(0.45250,0.3199602) (0.45750,0.3206089) (0.46250,0.3212072) (0.46750,0.3217634) (0.47250,0.3222462)
(0.47750,0.3227090) (0.48250,0.3230957) (0.48750,0.3234309) (0.49250,0.3237206) (0.49750,0.3239241)
(0.50250,0.3240780) (0.50750,0.3241964) (0.51250,0.3242665) (0.51750,0.3243485) (0.52250,0.3244863)
(0.52750,0.3246839) (0.53250,0.3250532) (0.53750,0.3255882) (0.54250,0.3263528) (0.54750,0.3274406)
(0.55250,0.3288437) (0.55750,0.3306364) (0.56250,0.3328659) (0.56750,0.3355137) (0.57250,0.3386367)
(0.57750,0.3422789) (0.58250,0.3463720) (0.58750,0.3509664) (0.59250,0.3560861) (0.59750,0.3616548)
(0.60250,0.3677388) (0.60750,0.3742661) (0.61250,0.3812681) (0.61750,0.3887937) (0.62250,0.3967364)
(0.62750,0.4051694) (0.63250,0.4141655) (0.63750,0.4236385) (0.64250,0.4336884) (0.64750,0.4443638)
(0.65250,0.4556850) (0.65750,0.4678173) (0.66250,0.4807112) (0.66750,0.4944767) (0.67250,0.5093652)
(0.67750,0.5252350) (0.68250,0.5422887) (0.68750,0.5606261) (0.69250,0.5803079) (0.69750,0.6014495)
(0.70250,0.6241806) (0.70750,0.6484471) (0.71250,0.6743478) (0.71750,0.7021306) (0.72250,0.7314108)
(0.72750,0.7623438) (0.73250,0.7948340) (0.73750,0.8288279) (0.74250,0.8642964) (0.74750,0.9010127)
(0.75250,0.9389422) (0.75750,0.9784038) (0.76250,1.0186922) (0.76750,1.0600660) (0.77250,1.1025110)
(0.77750,1.1460566) (0.78250,1.1908695) (0.78750,1.2369838) (0.79250,1.2846060) (0.79750,1.3345180)
(0.80250,1.3862315) (0.80750,1.4404566) (0.81250,1.4975760) (0.81750,1.5580087) (0.82250,1.6222305)
(0.82750,1.6906770) (0.83250,1.7636954) (0.83750,1.8421414) (0.84250,1.9260296) (0.84750,2.0150960)
(0.85250,2.1097353) (0.85750,2.2097632) (0.86250,2.3149849) (0.86750,2.4249688) (0.87250,2.5391752)
(0.87750,2.6570141) (0.88250,2.7783110) (0.88750,2.9016249) (0.89250,3.0267019) (0.89750,3.1530556)
(0.90250,3.2801841) (0.90750,3.4079522) (0.91250,3.5362259) (0.91750,3.6650249) (0.92250,3.7949766)
(0.92750,3.9255664) (0.93250,4.0572159) (0.93750,4.1899336) (0.94250,4.3232626) (0.94750,4.4566951)
(0.95250,4.5890965) (0.95750,4.7188632) (0.96250,4.8443614) (0.96750,4.9627646) (0.97250,5.0715382)
(0.97750,5.1679548) (0.98250,5.2490695) (0.98750,5.3124875) (0.99250,5.3560413) (0.99750,5.3781056)
}; \addlegendentry{\(N=96\)}
\addplot[rust, dashed,line width=1.1pt,no marks] coordinates {
(0.00250,0.2792377) (0.00750,0.2843011) (0.01250,0.2942930) (0.01750,0.3089483) (0.02250,0.3278806)
(0.02750,0.3506124) (0.03250,0.3765612) (0.03750,0.4050914) (0.04250,0.4355561) (0.04750,0.4672576)
(0.05250,0.4995316) (0.05750,0.5317396) (0.06250,0.5632916) (0.06750,0.5936634) (0.07250,0.6223950)
(0.07750,0.6491062) (0.08250,0.6734897) (0.08750,0.6953101) (0.09250,0.7143988) (0.09750,0.7306355)
(0.10250,0.7439541) (0.10750,0.7543559) (0.11250,0.7618500) (0.11750,0.7664907) (0.12250,0.7683609)
(0.12750,0.7675711) (0.13250,0.7642583) (0.13750,0.7585794) (0.14250,0.7507081) (0.14750,0.7408308)
(0.15250,0.7291471) (0.15750,0.7158585) (0.16250,0.7011702) (0.16750,0.6852843) (0.17250,0.6683971)
(0.17750,0.6507008) (0.18250,0.6323744) (0.18750,0.6135909) (0.19250,0.5945089) (0.19750,0.5752793)
(0.20250,0.5560420) (0.20750,0.5369259) (0.21250,0.5180398) (0.21750,0.4994873) (0.22250,0.4813490)
(0.22750,0.4636883) (0.23250,0.4465764) (0.23750,0.4300426) (0.24250,0.4141136) (0.24750,0.3988044)
(0.25250,0.3841244) (0.25750,0.3700762) (0.26250,0.3566576) (0.26750,0.3438640) (0.27250,0.3316926)
(0.27750,0.3201367) (0.28250,0.3091875) (0.28750,0.2988397) (0.29250,0.2890704) (0.29750,0.2798827)
(0.30250,0.2712405) (0.30750,0.2631166) (0.31250,0.2554804) (0.31750,0.2482899) (0.32250,0.2415066)
(0.32750,0.2350863) (0.33250,0.2289901) (0.33750,0.2231750) (0.34250,0.2176103) (0.34750,0.2122618)
(0.35250,0.2071061) (0.35750,0.2021425) (0.36250,0.1973586) (0.36750,0.1927559) (0.37250,0.1883473)
(0.37750,0.1841461) (0.38250,0.1801708) (0.38750,0.1764405) (0.39250,0.1729760) (0.39750,0.1697954)
(0.40250,0.1669130) (0.40750,0.1643321) (0.41250,0.1620555) (0.41750,0.1600766) (0.42250,0.1583949)
(0.42750,0.1569923) (0.43250,0.1558517) (0.43750,0.1549564) (0.44250,0.1542890) (0.44750,0.1538348)
(0.45250,0.1535737) (0.45750,0.1534917) (0.46250,0.1535764) (0.46750,0.1538196) (0.47250,0.1542079)
(0.47750,0.1547331) (0.48250,0.1553914) (0.48750,0.1561776) (0.49250,0.1570882) (0.49750,0.1581224)
(0.50250,0.1592808) (0.50750,0.1605687) (0.51250,0.1619912) (0.51750,0.1635572) (0.52250,0.1652761)
(0.52750,0.1671719) (0.53250,0.1692462) (0.53750,0.1715235) (0.54250,0.1740198) (0.54750,0.1767620)
(0.55250,0.1797719) (0.55750,0.1830759) (0.56250,0.1867050) (0.56750,0.1906897) (0.57250,0.1950628)
(0.57750,0.1998607) (0.58250,0.2051187) (0.58750,0.2108712) (0.59250,0.2171690) (0.59750,0.2240227)
(0.60250,0.2314711) (0.60750,0.2395518) (0.61250,0.2482907) (0.61750,0.2577175) (0.62250,0.2678678)
(0.62750,0.2787708) (0.63250,0.2904725) (0.63750,0.3030188) (0.64250,0.3164532) (0.64750,0.3308369)
(0.65250,0.3462913) (0.65750,0.3628344) (0.66250,0.3806093) (0.66750,0.3997049) (0.67250,0.4202657)
(0.67750,0.4424249) (0.68250,0.4663428) (0.68750,0.4921783) (0.69250,0.5201124) (0.69750,0.5503281)
(0.70250,0.5829909) (0.70750,0.6182666) (0.71250,0.6563352) (0.71750,0.6973675) (0.72250,0.7414639)
(0.72750,0.7887542) (0.73250,0.8393039) (0.73750,0.8932237) (0.74250,0.9505881) (0.74750,1.0113837)
(0.75250,1.0757405) (0.75750,1.1436784) (0.76250,1.2151172) (0.76750,1.2901946) (0.77250,1.3689518)
(0.77750,1.4513451) (0.78250,1.5374950) (0.78750,1.6273250) (0.79250,1.7210154) (0.79750,1.8187096)
(0.80250,1.9203335) (0.80750,2.0260670) (0.81250,2.1360484) (0.81750,2.2501329) (0.82250,2.3683200)
(0.82750,2.4904620) (0.83250,2.6161604) (0.83750,2.7449131) (0.84250,2.8759769) (0.84750,3.0083699)
(0.85250,3.1409011) (0.85750,3.2720619) (0.86250,3.4001859) (0.86750,3.5235218) (0.87250,3.6399150)
(0.87750,3.7473096) (0.88250,3.8437255) (0.88750,3.9269375) (0.89250,3.9949786) (0.89750,4.0461854)
(0.90250,4.0788880) (0.90750,4.0919135) (0.91250,4.0843084) (0.91750,4.0555955) (0.92250,4.0059897)
(0.92750,3.9360147) (0.93250,3.8468393) (0.93750,3.7403909) (0.94250,3.6190825) (0.94750,3.4859704)
(0.95250,3.3447086) (0.95750,3.1993268) (0.96250,3.0542451) (0.96750,2.9140169) (0.97250,2.7832380)
(0.97750,2.6663126) (0.98250,2.5673620) (0.98750,2.4898251) (0.99250,2.4365158) (0.99750,2.4093638)
}; \addlegendentry{\(N=288\)}
\addplot[black!55, dashdotted,line width=1.1pt,no marks] coordinates {
(0.00250,0.8711919) (0.00750,0.8793693) (0.01250,0.8953061) (0.01750,0.9181834) (0.02250,0.9468346)
(0.02750,0.9798172) (0.03250,1.0154913) (0.03750,1.0521056) (0.04250,1.0879040) (0.04750,1.1211950)
(0.05250,1.1504575) (0.05750,1.1743826) (0.06250,1.1919384) (0.06750,1.2023835) (0.07250,1.2052937)
(0.07750,1.2005483) (0.08250,1.1883087) (0.08750,1.1689578) (0.09250,1.1431123) (0.09750,1.1115290)
(0.10250,1.0750623) (0.10750,1.0346275) (0.11250,0.9911466) (0.11750,0.9455194) (0.12250,0.8985898)
(0.12750,0.8511049) (0.13250,0.8037378) (0.13750,0.7570498) (0.14250,0.7115009) (0.14750,0.6674481)
(0.15250,0.6251540) (0.15750,0.5848000) (0.16250,0.5464954) (0.16750,0.5102910) (0.17250,0.4761895)
(0.17750,0.4441569) (0.18250,0.4141335) (0.18750,0.3860399) (0.19250,0.3597880) (0.19750,0.3352855)
(0.20250,0.3124347) (0.20750,0.2911469) (0.21250,0.2713179) (0.21750,0.2528680) (0.22250,0.2357159)
(0.22750,0.2197800) (0.23250,0.2049866) (0.23750,0.1912605) (0.24250,0.1785356) (0.24750,0.1667530)
(0.25250,0.1558308) (0.25750,0.1457162) (0.26250,0.1363533) (0.26750,0.1276851) (0.27250,0.1196523)
(0.27750,0.1122096) (0.28250,0.1053074) (0.28750,0.0989061) (0.29250,0.0929596) (0.29750,0.0874377)
(0.30250,0.0823077) (0.30750,0.0775427) (0.31250,0.0731200) (0.31750,0.0690158) (0.32250,0.0652113)
(0.32750,0.0616893) (0.33250,0.0584378) (0.33750,0.0554377) (0.34250,0.0526778) (0.34750,0.0501449)
(0.35250,0.0478262) (0.35750,0.0457081) (0.36250,0.0437794) (0.36750,0.0420266) (0.37250,0.0404377)
(0.37750,0.0389985) (0.38250,0.0376986) (0.38750,0.0365277) (0.39250,0.0354746) (0.39750,0.0345302)
(0.40250,0.0336863) (0.40750,0.0329352) (0.41250,0.0322717) (0.41750,0.0316891) (0.42250,0.0311837)
(0.42750,0.0307523) (0.43250,0.0303917) (0.43750,0.0300999) (0.44250,0.0298746) (0.44750,0.0297141)
(0.45250,0.0296166) (0.45750,0.0295813) (0.46250,0.0296055) (0.46750,0.0296891) (0.47250,0.0298309)
(0.47750,0.0300294) (0.48250,0.0302846) (0.48750,0.0305958) (0.49250,0.0309632) (0.49750,0.0313880)
(0.50250,0.0318713) (0.50750,0.0324149) (0.51250,0.0330233) (0.51750,0.0337002) (0.52250,0.0344507)
(0.52750,0.0352808) (0.53250,0.0361980) (0.53750,0.0372110) (0.54250,0.0383284) (0.54750,0.0395601)
(0.55250,0.0409194) (0.55750,0.0424190) (0.56250,0.0440728) (0.56750,0.0458969) (0.57250,0.0479087)
(0.57750,0.0501286) (0.58250,0.0525795) (0.58750,0.0552820) (0.59250,0.0582684) (0.59750,0.0615694)
(0.60250,0.0652212) (0.60750,0.0692649) (0.61250,0.0737492) (0.61750,0.0787218) (0.62250,0.0842449)
(0.62750,0.0903804) (0.63250,0.0971977) (0.63750,0.1047856) (0.64250,0.1132283) (0.64750,0.1226314)
(0.65250,0.1330918) (0.65750,0.1447440) (0.66250,0.1577192) (0.66750,0.1721763) (0.67250,0.1882530)
(0.67750,0.2061753) (0.68250,0.2261460) (0.68750,0.2483990) (0.69250,0.2732001) (0.69750,0.3008607)
(0.70250,0.3317053) (0.70750,0.3661152) (0.71250,0.4044667) (0.71750,0.4472680) (0.72250,0.4950086)
(0.72750,0.5482033) (0.73250,0.6074582) (0.73750,0.6733576) (0.74250,0.7465391) (0.74750,0.8276505)
(0.75250,0.9173375) (0.75750,1.0161929) (0.76250,1.1248899) (0.76750,1.2439570) (0.77250,1.3738701)
(0.77750,1.5150093) (0.78250,1.6676283) (0.78750,1.8318080) (0.79250,2.0074472) (0.79750,2.1941783)
(0.80250,2.3914040) (0.80750,2.5981873) (0.81250,2.8132720) (0.81750,3.0350288) (0.82250,3.2614416)
(0.82750,3.4900973) (0.83250,3.7181885) (0.83750,3.9425307) (0.84250,4.1596028) (0.84750,4.3656019)
(0.85250,4.5565243) (0.85750,4.7282666) (0.86250,4.8767489) (0.86750,4.9980542) (0.87250,5.0885801)
(0.87750,5.1452034) (0.88250,5.1654065) (0.88750,5.1474740) (0.89250,5.0905332) (0.89750,4.9947092)
(0.90250,4.8611279) (0.90750,4.6919414) (0.91250,4.4902731) (0.91750,4.2601099) (0.92250,4.0061480)
(0.92750,3.7337707) (0.93250,3.4486570) (0.93750,3.1566991) (0.94250,2.8637952) (0.94750,2.5756603)
(0.95250,2.2976407) (0.95750,2.0346737) (0.96250,1.7910894) (0.96750,1.5706310) (0.97250,1.3764987)
(0.97750,1.2112372) (0.98250,1.0768471) (0.98750,0.9748305) (0.99250,0.9062569) (0.99750,0.8718282)
}; \addlegendentry{\(N=768\)}
\addplot[teal, solid,line width=1.1pt,no marks] coordinates {
(0.00250,1.1568929) (0.00750,1.1649932) (0.01250,1.1806968) (0.01750,1.2030458) (0.02250,1.2306767)
(0.02750,1.2619120) (0.03250,1.2948517) (0.03750,1.3274931) (0.04250,1.3578393) (0.04750,1.3840129)
(0.05250,1.4043383) (0.05750,1.4174242) (0.06250,1.4222136) (0.06750,1.4180193) (0.07250,1.4045156)
(0.07750,1.3817379) (0.08250,1.3500509) (0.08750,1.3101126) (0.09250,1.2627833) (0.09750,1.2091162)
(0.10250,1.1502593) (0.10750,1.0874377) (0.11250,1.0218529) (0.11750,0.9546858) (0.12250,0.8870308)
(0.12750,0.8198850) (0.13250,0.7541119) (0.13750,0.6904429) (0.14250,0.6294764) (0.14750,0.5716714)
(0.15250,0.5173517) (0.15750,0.4667286) (0.16250,0.4199028) (0.16750,0.3768863) (0.17250,0.3376149)
(0.17750,0.3019628) (0.18250,0.2697618) (0.18750,0.2408040) (0.19250,0.2148665) (0.19750,0.1917139)
(0.20250,0.1711014) (0.20750,0.1527877) (0.21250,0.1365487) (0.21750,0.1221636) (0.22250,0.1094341)
(0.22750,0.0981612) (0.23250,0.0881894) (0.23750,0.0793556) (0.24250,0.0715311) (0.24750,0.0645828)
(0.25250,0.0584199) (0.25750,0.0529462) (0.26250,0.0480776) (0.26750,0.0437401) (0.27250,0.0398785)
(0.27750,0.0364375) (0.28250,0.0333691) (0.28750,0.0306305) (0.29250,0.0281886) (0.29750,0.0260078)
(0.30250,0.0240605) (0.30750,0.0223191) (0.31250,0.0207667) (0.31750,0.0193804) (0.32250,0.0181429)
(0.32750,0.0170381) (0.33250,0.0160533) (0.33750,0.0151764) (0.34250,0.0143950) (0.34750,0.0136990)
(0.35250,0.0130808) (0.35750,0.0125305) (0.36250,0.0120412) (0.36750,0.0116054) (0.37250,0.0112171)
(0.37750,0.0108702) (0.38250,0.0105591) (0.38750,0.0102796) (0.39250,0.0100279) (0.39750,0.0098004)
(0.40250,0.0095946) (0.40750,0.0094080) (0.41250,0.0092395) (0.41750,0.0090877) (0.42250,0.0089518)
(0.42750,0.0088314) (0.43250,0.0087262) (0.43750,0.0086359) (0.44250,0.0085607) (0.44750,0.0085002)
(0.45250,0.0084547) (0.45750,0.0084238) (0.46250,0.0084074) (0.46750,0.0084053) (0.47250,0.0084176)
(0.47750,0.0084439) (0.48250,0.0084842) (0.48750,0.0085386) (0.49250,0.0086072) (0.49750,0.0086902)
(0.50250,0.0087882) (0.50750,0.0089015) (0.51250,0.0090312) (0.51750,0.0091783) (0.52250,0.0093444)
(0.52750,0.0095309) (0.53250,0.0097402) (0.53750,0.0099750) (0.54250,0.0102382) (0.54750,0.0105335)
(0.55250,0.0108659) (0.55750,0.0112403) (0.56250,0.0116631) (0.56750,0.0121406) (0.57250,0.0126814)
(0.57750,0.0132937) (0.58250,0.0139885) (0.58750,0.0147761) (0.59250,0.0156699) (0.59750,0.0166840)
(0.60250,0.0178372) (0.60750,0.0191468) (0.61250,0.0206347) (0.61750,0.0223289) (0.62250,0.0242639)
(0.62750,0.0264766) (0.63250,0.0290112) (0.63750,0.0319271) (0.64250,0.0352985) (0.64750,0.0391986)
(0.65250,0.0437370) (0.65750,0.0490446) (0.66250,0.0552627) (0.66750,0.0625719) (0.67250,0.0711714)
(0.67750,0.0813224) (0.68250,0.0933478) (0.68750,0.1075448) (0.69250,0.1243368) (0.69750,0.1441797)
(0.70250,0.1675898) (0.70750,0.1951555) (0.71250,0.2275364) (0.71750,0.2654794) (0.72250,0.3098172)
(0.72750,0.3613851) (0.73250,0.4211554) (0.73750,0.4901275) (0.74250,0.5693545) (0.74750,0.6598527)
(0.75250,0.7626976) (0.75750,0.8788908) (0.76250,1.0093719) (0.76750,1.1549379) (0.77250,1.3162188)
(0.77750,1.4936378) (0.78250,1.6873017) (0.78750,1.8970006) (0.79250,2.1221293) (0.79750,2.3616316)
(0.80250,2.6139662) (0.80750,2.8770866) (0.81250,3.1484158) (0.81750,3.4248545) (0.82250,3.7028061)
(0.82750,3.9782230) (0.83250,4.2466766) (0.83750,4.5034500) (0.84250,4.7436519) (0.84750,4.9623494)
(0.85250,5.1547172) (0.85750,5.3161949) (0.86250,5.4426512) (0.86750,5.5305442) (0.87250,5.5770713)
(0.87750,5.5803006) (0.88250,5.5392855) (0.88750,5.4541235) (0.89250,5.3259963) (0.89750,5.1571686)
(0.90250,4.9509365) (0.90750,4.7114910) (0.91250,4.4438429) (0.91750,4.1536011) (0.92250,3.8467803)
(0.92750,3.5296173) (0.93250,3.2083309) (0.93750,2.8889764) (0.94250,2.5771539) (0.94750,2.2779767)
(0.95250,1.9959176) (0.95750,1.7347946) (0.96250,1.4976299) (0.96750,1.2868128) (0.97250,1.1041195)
(0.97750,0.9507590) (0.98250,0.8275118) (0.98750,0.7348397) (0.99250,0.6729700) (0.99750,0.6420481)
}; \addlegendentry{\(N=1400\)}

\end{axis}
\end{tikzpicture}
\caption{Равновесная плотность химического порядка после снятия зонтика
при \(W<0{,}40\) и свободном \(c/a\). Каждая кривая построена при своём
поле сосуществования по весам на один вариант порядка и нормирована на
единицу. С ростом \(N\) уменьшается вероятность промежуточного порядка.
При определении поля сравниваются площади под двумя пиками, а не их высоты.}
\label{fig:order}
\end{figure}

\subsection{Проверки положения границ}

Два независимых набора зонтичных траекторий на 1400 атомах при
$W<0{,}40$ дали для $\alpha$ значения $14{,}89$ и $15{,}08$
ат.\%, а для $\alpha_2$ --- $22{,}37$ и $22{,}53$ ат.\%.
При переносе разделяющей поверхности от $Q=0{,}34$ до $0{,}54$
каждая граница меняется максимум на $0{,}02$ п.п. Если в той же
ячейке потребовать равенства \emph{суммарных} весов бассейнов,
получаются $14{,}90$ и $22{,}34$ ат.\%; основной расчёт учитывает
вырождение четырёх вариантов порядка в $\alpha_2$.

Две проверки использовали распределение состава без деления по $Q$.
Пересечение зависимостей $\langle x\rangle(\Delta\mu)$ для
$N=768$ и 1400 при $W<0{,}40$ дало
$3{,}29397$ эВ с интервалом $3{,}29235$--$3{,}29640$ эВ.
Вторая проверка использует четвёртую кумулянту состава
$\kappa_4=\langle(x-\langle x\rangle)^4\rangle-
3\,\mathrm{Var}(x)^2$ в максимуме дисперсии. Для равновесной смеси
двух одинаково взвешенных узких пиков расстояние между ними
оценивается как $(-8\kappa_4)^{1/4}$. Эта оценка дала
$7{,}48$ и $7{,}46$ п.п. при $N=768$ и 1400 соответственно.
Она подтверждает ширину по разности фазовых средних, $7{,}46$ п.п.,
с помощью другого наблюдаемого на тех же выборках.

В независимом тесте для каждой величины $W$ запущено по четыре
случайных и четыре упорядоченных состояния на 1400 атомах.
По химическому порядку они эволюционировали свободно. Для каждой
фазы состав оценивался отдельно с поправкой на разницу между
заданным полем и полем сосуществования из зонтичных окон.
При $W<0{,}40$ получены $14{,}987$ и $22{,}489$ ат.\% Al,
с 95\%-ми интервалами по репликам и полю $14{,}88$--$15{,}11$
и $22{,}29$--$22{,}64$ ат.\%. При $W<0{,}35$ получены
$14{,}914$ и $22{,}486$ ат.\%; две первоначально упорядоченные
реплики перешли в $\alpha$ во время термализации. В производственной
части переходов между ядрами бассейнов не обнаружено.
Эти траектории проверяют \emph{составы фаз при найденном поле};
само поле определено из относительных статистических весов зонтичных
окон.

\subsection{Сопоставление с C46 при согласованной геометрии}
\label{sec:bulk-c46}

Для C46 повторно восстановлена общая касательная по 13 авторским
окнам при 1000 К, $p=0$ и $N=5120$. Составы совпали с исходным
результатом до точности записи. Геометрический контроль прямого
расчёта сохранял отношение осей C46, $c/a=4{,}6046/2{,}9446
=1{,}56374$, при свободном объёме. Для этого контроля поле
сосуществования уточнялось независимо, включая две серии окон
непосредственно вблизи $3{,}322$ эВ.

\begin{table}[htbp]
\centering\small
\caption{Геометрия ячейки и составы. Для MTP здесь $N=1400$,
$W<0{,}40$. В квадратных скобках приведены условные 95\%-е
статистические интервалы. Ширина $\Delta x=x_{\alpha_2}-x_{\alpha}$
дана в п.п.}
\label{tab:bulk-c46}
\begin{tabular}{@{}l c c r@{}}
\toprule
Метод и форма & $\alpha$, ат.\% Al & $\alpha_2$, ат.\% Al & $\Delta x$\\
\midrule
MTP, свободное $c/a$ & $14{,}98\ [14{,}92;\,15{,}06]$ & $22{,}44\ [22{,}33;\,22{,}55]$ & $7{,}46$\\
MTP, $c/a$ как в C46 & $16{,}08\ [16{,}03;\,16{,}14]$ & $22{,}16\ [22{,}06;\,22{,}24]$ & $6{,}08$\\
C46, заданное $c/a$ & $14{,}96\ [14{,}93;\,14{,}99]$ & $21{,}96\ [21{,}91;\,22{,}01]$ & $7{,}00$\\
\bottomrule
\end{tabular}
\end{table}

В фиксированной геометрии увеличение ячейки с 768 до 1400 атомов
сдвинуло оценки на $-0{,}036$ и $-0{,}065$ п.п. Использование
только двух серий вблизи поля сосуществования дало
$16{,}091$ и $22{,}169$ ат.\% при $N=1400$. Это согласуется
с объединёнными оценками $16{,}083$ и $22{,}159$ ат.\%.
Поле в конфигурационной нормировке равно $3{,}32224$ эВ.

Сравнение с согласованным $c/a$ даёт совокупный сдвиг от C46
к прямому MTP $+1{,}12$ п.п. для $\alpha$ и $+0{,}20$ п.п.
для $\alpha_2$. При последующей релаксации отношения осей в прямом
расчёте граница $\alpha$ сдвигается на $-1{,}10$ п.п.,
а граница $\alpha_2$ --- на $+0{,}29$ п.п. Их 95\%-е интервалы
составляют от $-1{,}19$ до $-1{,}00$ и от $+0{,}14$ до
$+0{,}43$ п.п. соответственно.
В итоге относительно C46 получаем $+0{,}02$ и $+0{,}49$ п.п.
Близость двух конечных оценок для $\alpha$ возникает вследствие
компенсации сдвигов противоположного знака. Ширина области
сосуществования при освобождении осей возрастает на $1{,}39$ п.п.
с интервалом $1{,}22$--$1{,}54$ п.п.

У обоих расчётов один исходный MTP. При совпадающем $c/a$ остаются
различия в представлении химической энергии и колебаний
кластерным разложением C46, в квазигармоническом приближении и
в статистике ядер (квантовая в C46, классическая в прямом
расчёте). Измеренная разность характеризует их совокупный вклад.

\paragraph{Прослеживаемость.}
На UMT под учётной записью \texttt{u3129} завершены 26 производственных
серий и независимых контролей. Исходные идентификаторы Slurm,
SHA-256 архивов и число проходов по сериям сохранены в
\texttt{run\_provenance.json}. Числа для таблиц получены из
\texttt{summary\_statistics.json} и \texttt{c46\_reference.json},
проверки по составу и прямым траекториям ---
из \texttt{phase\_checks.json} и \texttt{endpoint\_controls.json}.
Вычислитель энергии отдельно сопоставлен с MLIP3 на 123 конфигурациях
из 8, 128 и 432 атомов, включая смены типа, смещения и изменения
объёма. Максимальное абсолютное расхождение составило
$1{,}36\cdot10^{-12}$ эВ на всю ячейку.

\section{Независимые методы и дальнейшие расчёты}
\label{sec:outlook}

\subsection{Какие методы применимы к сплаву}

Для сплава замещения абсолютная свободная энергия должна учитывать и движение
атомов, и перестановки Ti и Al. При фиксированных составе и форме ячейки
статистическая сумма складывается по химическим конфигурациям; каждая
конфигурация входит со своим позиционным интегралом. Кристалл Эйнштейна в
методе Френкеля--Ладда даёт аналитическую опору для позиционной части.
Если атомы навсегда закреплены за выбранной химической конфигурацией, такой
расчёт определяет её свободную энергию. Для равновесного сплава дополнительно
требуется взвесить конфигурации со всеми разрешёнными антиструктурными
дефектами. Примеры совместного учёта химических и позиционных степеней
свободы приведены в работах Ли и Скандоло для неидеальных растворов и
Менона и соавторов для Au--Cu~\cite{frenkel1984,li2024,menon2026}.

Для независимой проверки Ti--Al подходит следующий обратимый путь.
В начале все атомы взаимодействуют как Ti, но заданное число \(k\) узлов
несёт метку Al. В конце каждая метка задаёт настоящий тип атома в MTP.
При симметричных обменах меток, одинаковых фиктивных массах, фиксированных
\(T,V\) и общем ограничении на смещения разность свободных энергий равна
\begin{align}
 U_\lambda &=(1-\lambda)U_{\rm Ti}+\lambda U_{\rm MTP}, \notag\\
 F_{\rm MTP}(k)-F_{\rm Ti}
   &=-k_{\rm B}T\ln {N\choose k}
      +\int_0^1 \big\langle U_{\rm MTP}-U_{\rm Ti}
                        \big\rangle_{\lambda,k}\,d\lambda .
 \label{eq:alloy-alchemical}
\end{align}
Все энергии здесь относятся к целой ячейке. Биномиальный множитель следует
из числа равновероятных размещений меток в начальном состоянии. Если уже в
начальной точке введено дополнительное поле, выбирающее один вариант
упорядочения, для него надо вычислить собственную сумму по заполнениям и
учесть работу снятия поля. Перестановки и смещения при каждом \(\lambda\)
принимают с энергией \(U_\lambda\). Устойчивость выбранного
ГПУ-эталона на пути проверяют по структуре и по согласованности прямого
и обратного проходов. Для свободной ячейки можно
интегрировать при постоянных \(T,p\) и получить разность энергий Гиббса.
Геометрические ограничения и обработку глобального переноса следует
согласовать между началом и концом пути. Применение такого пути к Au--Cu
описано в препринте~\cite{menon2026}.

Есть две альтернативы для других постановок. В сплаве Fe--Cr Садиг и Эрхарт
ограничивали флуктуации состава в полубольшом каноническом ансамбле и
интегрировали химический потенциал. Это позволяет измерять свободную энергию
в том числе у ячеек, содержащих межфазную границу. Для получения объёмной
бинодали становятся существенными площадь границы и энергия когерентных
напряжений~\cite{sadigh2012}. Если требуется искать также конкурирующие
кристаллические структуры, можно применить вложенную выборку
(\emph{nested sampling}). Розенброк и соавторы построили таким способом
диаграмму Ag--Pd с MTP. Для применения к Ti--Al понадобились бы проверка
потенциала на всех затронутых структурах и отдельное исследование
сходимости структурного поиска~\cite{rosenbrock2021}.

Существует опубликованное сравнение на самом Ti--Al. Николс и соавторы
привели для MTP при 1000~K значения 15,2 и 23,4~ат.\% Al
для направлений \(\alpha\to\alpha_2\) и \(\alpha_2\to\alpha\).
Авторы брали последний наблюдаемый состав перед скачком внутри петли
гистерезиса в ячейке из 1024 атомов. На эти значения влияют время
наблюдения и зарождение новой фазы. В настоящей работе поле определено
по равенству статистических весов на один вариант упорядочения.
Для сопоставления работ следует сохранить оба критерия и указать метод
рядом с каждым числом~\cite{nichols2026}.

\subsection{Последовательность контрольных экспериментов}

\paragraph{1. Разложить расхождение с C46 при 1000~K.}
Начать с одинаковых \(p=0\) и фиксированного \(c/a=1{,}56374\).
В этой геометрии для границы \(\alpha\) остаются примерно \(1{,}12\)
процентного пункта, тогда как эффект освобождения \(c/a\) уже измерен.
Из исходных фононных частот тех же обучающих структур и объёмов получить
два набора меток. Для устойчивой моды частоты \(\omega\) они равны
\[
 f_{\rm cl}=k_{\rm B}T\ln(\beta\hbar\omega),\qquad
 f_{\rm q}=k_{\rm B}T\ln\!\left[2\sinh(\beta\hbar\omega/2)\right],
 \qquad \beta=(k_{\rm B}T)^{-1}.
\]
Сохранить набор кластеров, веса регрессии, обработку нулевых мод и
объёмную сетку. Повторное определение обеих границ с классическими
метками покажет, как замена статистики ядер меняет предсказание C46 при
неизменной схеме подгонки. Затем на независимых декорациях из обеих фаз
сопоставить предсказание классического CE с явно вычисленным гармоническим
вкладом того же MTP. В малой ячейке можно выполнить полную сумму по
заполнениям. Для большей ячейки коррекцию энергии по декорациям можно
перевзвесить при достаточном перекрытии выборок. Разность свободных
энергий, её производная по составу и сдвиг обеих границ дадут оценку
погрешности кластерной аппроксимации.

Для оставшегося ангармонического вклада сравнить гармонический эталон
с прямым классическим MTP в \emph{одной и той же} области \(W\).
Для устойчивой декорации с гауссовским эталоном переход к этой области
добавляет \(-k_{\rm B}T\ln P_{\rm harm}(W<W_{\max})\) к её свободной
энергии. Эту поправку следует отслеживать отдельно от изменения силового
закона. Декорации с неустойчивыми модами требуют собственного
ангармонического эталона либо явного изменения области сопоставления.
В результате получим цепочку сравнений по обеим границам. Численные
вклады относятся к выбранному порядку замены моделей.

\paragraph{2. Замкнуть независимый цикл свободных энергий.}
Выбрать по нескольку составов по обе стороны от каждого контакта, провести
интегрирование по \(\lambda\) в обе стороны с обменом типов и контролем
порядка. Сопоставить разность свободных энергий между выбранными составами
с интегралом \(\int\Delta\mu\,dx\) по той же ветви и в тех же единицах
на атом. Если ветвь фиксируется дополнительным смещающим полем, учесть
его работу и при интегрировании, и при восстановлении исходных весов.
При резкой перестройке порядка внутри пути можно добавить окна по \(Q\).
Критерием согласия служит замыкание цикла в пределах суммарного интервала
ошибок по независимым траекториям. Так можно проверить общую нормировку
фаз и одновременно оценить практическую стоимость прямого TDI.

\paragraph{3. Определить поправку исходного MTP к DFT.}
Сформировать выборку термически декоррелированных конфигураций с сохранением
координат, ячеек и химических типов. Нужны состояния обеих фаз в окрестности
их равновесных составов. Для каждого состояния вычислить полную энергию
MTP и DFT на \emph{одинаковых} координатах; исходной целью DFT выбрать
функционал и энергетическую конвенцию обучающего набора потенциала
Qi~\cite{qi2023}. Начать с ячеек, доступных для DFT, и проверить
переносимость поправки по размеру, составу и форме ячейки. Для фиксированных
состава, фазы и механического ансамбля формула возмущений даёт
\[
 \delta G_i=-k_{\rm B}T
   \ln\left\langle \exp[-\beta(U_{\rm DFT}-U_{\rm MTP})]\right\rangle_{i,{\rm MTP}} .
\]
Здесь давление, область фазы, ограничения и мера одинаковы в обоих
ансамблях. При фиксированном объёме та же формула даёт \(\delta F_i\).
В экспоненте стоит разность энергий \emph{всей ячейки}. Оценить долю
веса крупнейших конфигураций, эффективный объём выборки и устойчивость
к разбиению её по независимым репликам. Если вес сосредоточен в
единичных кадрах, потребуются промежуточные гамильтонианы и
двусторонняя оценка. Подобная привязка свободных энергий к DFT была
продемонстрирована для сплава TaVCrW~\cite{zhu2024}.

После построения \(\delta g_i(x)=\delta G_i(x)/N\) следует снова найти
общую касательную. Для масштаба задачи малая \emph{разность} поправок
между фазами в 1~мэВ/атом при измеренном разрыве состава
\(\Delta x\simeq0{,}0746\) сдвигает поле сосуществования примерно на
\(0{,}001/0{,}0746\simeq0{,}0134\)~эВ на замену. Это оценка
чувствительности. Сдвиг самих составов зависит ещё от производных
\(\delta g_i(x)\) и от кривизны каждой ветви.

\paragraph{4. Расширить область проверки и сопоставить с измерениями.}
Короткодействующую часть потенциала стоит проверить отдельными DFT-точками
на траектории сближения Ti--Ti. Для варианта с исправленным отталкиванием
надо заново измерить границы и их зависимость от \(W\), сохранив результаты
исходного MTP как отдельный набор. Для температурного хода сольвуса
целесообразно сначала получить независимые опорные точки по обе стороны
от 1000~K. Между ними можно интегрировать уравнение Гиббса--Дюгема
при постоянном давлении~\cite{nam2007},
\[
 \frac{d(\beta\Delta\mu)}{d\beta}
     =\frac{h_{\alpha_2}-h_\alpha}{x_{\alpha_2}-x_\alpha}.
\]
Здесь \(h_i\) и \(\Delta\mu\) используют согласованную термодинамическую
конвенцию и гамильтониан без явной зависимости от температуры. При слиянии
составов правую часть труднее оценить, и требуется новый прямой расчёт.
С экспериментом следует сопоставлять обе ветви при совпадающих температурах
и в атомных процентах с указанием состава примесей, времени отжига и
погрешности определения фазовых составов. Работа Онумы и соавторов
содержит первичные измерения на Ti--Al с низким содержанием кислорода и
отдельную термодинамическую оценку~\cite{ohnuma2000}. После привязки к DFT
останутся вопросы выбора функционала, электронного температурного вклада
и условий эксперимента. Они задают предел интерпретации сравнения с опытом.

\section*{Расчётные данные и воспроизведение}
\label{sec:archive}
Данные организованы относительно каталога проекта
\nolinkurl{phase_reproduction/}. Исходный файл потенциала —
\nolinkurl{inputs/fitted.mtp}, тип 0 соответствует Ti, тип 1 — Al.
Его SHA-256 по реестру запусков и независимой проверке энергии равен
\begin{center}
{\small\ttfamily 704f7fd900c7446286066b741084c30ebc7fcf26e5b9b8fd2f29a18ba4ffd661}.
\end{center}
\begin{itemize}
\item Малые ячейки. \nolinkurl{alloy_ti/targeted_comparison.json} содержит
оценки с дополнительной выборкой для четырёх наиболее весомых классов;
\nolinkurl{alloy_ti/cage_sensitivity.json} — независимые оценки при двух
ограничениях. Геометрию короткой связи и независимый пересчёт энергии
сохраняет \nolinkurl{cluster/runs/alloy-ti-geometry-01/outputs/}.
\item Атомистическая бинодаль. В \nolinkurl{solvus_1000/series/} каждой
серии соответствует каталог с \texttt{windows.csv}, записями измерений и
конечными состояниями. В \nolinkurl{solvus_1000/run_provenance.json}
проверяются номера заданий UMT, хеши полученных архивов и состав серий.
В \nolinkurl{solvus_1000/summary_statistics.json} приведены результаты
для десяти ансамблей и попарные изменения границ.
\item Сравнение с C46. \nolinkurl{solvus_1000/c46_reference.json}
содержит тринадцать исходных пар «состав — химический потенциал»,
результат общей касательной, условные интервалы и хеши авторских файлов.
\end{itemize}

Для повторения основного анализа нужны Python, NumPy и SciPy. Следующие
команды PowerShell запускаются из \nolinkurl{phase_reproduction/solvus_1000/}.
Они восстанавливают обе изотермы для 1400 атомов со свободным \(c/a\) и
их условные статистические интервалы из завершённых траекторий.
{\small
\begin{verbatim}
$m = 's1400w35a','s1400w35b','s1400w40a','s1400w40b' |
     ForEach-Object { "series/$_/windows.csv" }
python analysis.py @m --bootstrap 192 --seed 26713 --output replay.json
\end{verbatim}
}

\begin{thebibliography}{99}

\bibitem{qi2023}
J.~Qi, Z.~H.~Aitken, Q.~Pei et al.
Machine learning moment tensor potential for modeling dislocation and fracture
in \(\mathrm{L1}_0\)-TiAl and \(\mathrm{D0}_{19}\)-\(\mathrm{Ti}_3\mathrm{Al}\) alloys.
\emph{Physical Review Materials} \textbf{7}, 103602 (2023).
\href{https://doi.org/10.1103/PhysRevMaterials.7.103602}{DOI}.

\bibitem{mbar2008}
M.~R.~Shirts, J.~D.~Chodera.
Statistically optimal analysis of samples from multiple equilibrium states.
\emph{Journal of Chemical Physics} \textbf{129}, 124105 (2008).
\href{https://doi.org/10.1063/1.2978177}{DOI}.

\bibitem{janke1993}
W.~Janke.
Accurate first-order transition points from finite-size data without power-law corrections.
\emph{Physical Review B} \textbf{47}, 14757 (1993).
\href{https://doi.org/10.1103/PhysRevB.47.14757}{DOI}.

\bibitem{borgs1992}
C.~Borgs, W.~Janke.
New method to determine first-order transition points from finite-size data.
\emph{Physical Review Letters} \textbf{68}, 1738 (1992).
\href{https://doi.org/10.1103/PhysRevLett.68.1738}{DOI}.

\bibitem{frenkel1984}
D.~Frenkel, A.~J.~C.~Ladd.
New Monte Carlo method to compute the free energy of arbitrary solids.
Application to the fcc and hcp phases of hard spheres.
\emph{Journal of Chemical Physics} \textbf{81}, 3188 (1984).
\href{https://doi.org/10.1063/1.448024}{DOI}.

\bibitem{li2024}
Z.~Li, S.~Scandolo.
Efficient determination of free energies of non-ideal solid solutions via hybrid
Monte Carlo simulations.
\emph{Computer Physics Communications} \textbf{304}, 109307 (2024).
\href{https://doi.org/10.1016/j.cpc.2024.109307}{DOI}.

\bibitem{menon2026}
S.~Menon, M.~Poul, T.~Hickel, J.~Neugebauer, R.~Drautz.
Computing binary alloy phase diagrams with explicit configurational and
vibrational entropy.
Препринт, версия от 29 октября 2026~г.
\href{https://arxiv.org/abs/2607.14795v4}{arXiv:2607.14795v4}.

\bibitem{sadigh2012}
B.~Sadigh, P.~Erhart.
Calculation of excess free energies of precipitates via direct thermodynamic
integration across phase boundaries.
\emph{Physical Review B} \textbf{86}, 134204 (2012).
\href{https://doi.org/10.1103/PhysRevB.86.134204}{DOI}.

\bibitem{rosenbrock2021}
C.~W.~Rosenbrock, K.~Gubaev, A.~V.~Shapeev et al.
Machine-learned interatomic potentials for alloys and alloy phase diagrams.
\emph{npj Computational Materials} \textbf{7}, 24 (2021).
\href{https://doi.org/10.1038/s41524-020-00477-2}{DOI}.

\bibitem{nichols2026}
M.~Nichols, M.~S.~Nitol, S.~J.~Fensin, C.~D.~Barrett, D.~E.~Dickel.
Predicting the Ti-Al binary phase diagram with an artificial neural network potential.
\emph{Metals} \textbf{16}, 140 (2026), разд.~2, табл.~2.
\href{https://doi.org/10.3390/met16020140}{DOI}.

\bibitem{zhu2024}
L.-F.~Zhu, F.~Körmann, Q.~Chen, M.~Selleby, J.~Neugebauer, B.~Grabowski.
Accelerating ab initio melting property calculations with machine learning:
application to the high entropy alloy TaVCrW.
\emph{npj Computational Materials} \textbf{10}, 274 (2024).
\href{https://doi.org/10.1038/s41524-024-01464-7}{DOI}.

\bibitem{nam2007}
H.-S.~Nam, M.~I.~Mendelev, D.~J.~Srolovitz.
Solid-liquid phase diagrams for binary metallic alloys: Adjustable interatomic potentials.
\emph{Physical Review B} \textbf{75}, 014204 (2007).
\href{https://doi.org/10.1103/PhysRevB.75.014204}{DOI}.

\bibitem{ohnuma2000}
I.~Ohnuma, Y.~Fujita, H.~Mitsui, K.~Ishikawa, R.~Kainuma, K.~Ishida.
Phase equilibria in the Ti--Al binary system.
\emph{Acta Materialia} \textbf{48}, 3113--3123 (2000).
\href{https://doi.org/10.1016/S1359-6454(00)00118-X}{DOI}.

\end{thebibliography}

\end{document}

